%% Sections 4, 5, 6: Agent Workflow Patterns

\section{Agent Workflow --- Streaming Execution}\label{sec:streaming}

In a streaming execution model, the wrapper monitors the LLM's output token-by-token as it is generated. As soon as a complete tool call is detected in the stream, the wrapper immediately pauses the LLM's generation, executes the tool, appends the result to the conversation context, and resumes generation. This cycle repeats until the LLM emits a terminate signal indicating no further actions are required.

The primary advantage of streaming execution is responsiveness. Users see output appearing in real time rather than waiting for a complete response. Tools fire as soon as they are needed, minimizing latency between decision and action. This creates a more interactive experience where intermediate results can inform subsequent reasoning steps without waiting for the entire generation to complete.

However, streaming execution requires more sophisticated infrastructure. The wrapper must parse incomplete output, maintain state across pause-resume cycles, and handle potential errors in partially generated tool calls. The LLM must also be capable of resuming generation coherently after external results are injected mid-stream.

Figure~\ref{fig:streaming} illustrates this workflow. The LLM generates output continuously, monitored for tool calls. Upon detection, execution is handed off to the tool, results are appended, and generation resumes. The process terminates only when the LLM explicitly signals completion.

\section{Agent Workflow --- Batch Execution}\label{sec:batch}

Batch execution follows a simpler but less responsive pattern. The LLM generates its entire response in one pass until it naturally stops. Only then does the wrapper parse the complete output, extract all tool calls present, and execute them serially one after another. All tool results are collected and fed back to the LLM simultaneously, prompting a new generation cycle. This process repeats until the LLM produces output with no tool calls, indicating task completion.

This is the classic chatbot model familiar from interactive coding assistants: the agent writes code, attempts to run it, receives error messages, and tries again. Batch execution is significantly simpler to implement than streaming because it requires no real-time parsing or state management. The wrapper operates on complete, well-formed outputs and can use standard parsing techniques.

The trade-off is latency. Users must wait for the entire generation to complete before any tools execute, and all tools must finish before the next reasoning step begins. For complex tasks requiring multiple tool invocations, this can result in noticeable delays between observable progress. However, for tasks where tool calls are infrequent or where thoughtful planning precedes action, batch execution provides a robust and maintainable architecture.

Figure~\ref{fig:batch} shows the batch workflow. The LLM produces complete output, which is parsed to extract all tool calls. These are executed serially, producing a combined result set that is fed back to the LLM to trigger the next generation cycle.

\section{Agent Workflow --- Parallel Execution}\label{sec:parallel_exec}

Parallel execution extends the batch model by introducing concurrency. Each tool call generated by the LLM can carry a flag indicating whether it can execute in parallel with other operations or requires blocking synchronous execution. When the LLM generates a tool call marked for parallel execution, the wrapper launches it asynchronously and allows the LLM to continue generating additional output or tool calls. When a non-parallel tool call is encountered or the LLM emits a terminate signal, the wrapper stops, waits for all pending parallel operations to complete, executes any blocking tool, and feeds all accumulated results back to the LLM.

This architecture combines the responsiveness of streaming execution with the throughput advantages of concurrency. Multiple independent operations—such as reading several files, running tests, or querying external APIs—can proceed simultaneously without blocking each other or the LLM's reasoning process. Critical operations that depend on prior results or require exclusive access to resources can still be executed synchronously to maintain correctness.

The complexity lies in managing parallel task lifetimes, handling errors from concurrent operations, and ensuring deterministic behavior when multiple results arrive out of order. The wrapper must track which parallel operations are outstanding, aggregate their results, and present them coherently to the LLM. The LLM itself must understand which operations can safely run concurrently and which require sequential ordering.

Figure~\ref{fig:parallel_workflow} illustrates this model. Tools marked with a parallel flag execute concurrently while generation continues. A synchronization barrier is inserted before any blocking tool, ensuring all prior parallel operations complete before the blocking operation begins. All results are then returned to the LLM together, maintaining a consistent view of the execution state.

\begin{figure*}[t]
\centering
\begin{tikzpicture}[node distance=1.5cm and 1.2cm]

\node[agentbox] (llm1) {LLM generates...};
\node[flag, right=of llm1] (detect) {Tool call detected};
\node[agentbox=horseshoe, right=of detect] (exec) {Execute tool};
\node[agentbox=atlantic, right=of exec] (append) {Append result};

\node[agentbox=congaree, below=0.8cm of llm1] (term1) {Terminate};

\draw[dataarrow] (llm1) -- (detect);
\draw[dataarrow] (detect) -- (exec);
\draw[dataarrow] (exec) -- (append);
\draw[dataarrow] (append.north) -- ++(0,0.8) -| (llm1.north);
\draw[dataarrow, dashed] (llm1) -- (term1);

\end{tikzpicture}
\caption{Streaming execution workflow. The LLM generates output incrementally; as soon as a tool call is detected, execution pauses, the tool runs, results are appended, and generation resumes. The process repeats until the LLM emits a terminate signal.}
\label{fig:streaming}
\end{figure*}

\begin{figure*}[t]
\centering
\begin{tikzpicture}[node distance=0.8cm and 0.8cm]

\node[agentbox] (llm2) {LLM};
\node[agentbox=bk70, right=of llm2] (output) {Full output};
\node[agentbox=atlantic, right=of output] (parse) {Parse calls};
\node[pbar, right=of parse] (t1) {T1};
\node[pbar, right=0.4cm of t1] (t2) {T2};
\node[pbar, right=0.4cm of t2] (t3) {T3};
\node[agentbox=congaree, right=of t3] (allres) {All results};

% Terminate under LLM
\node[agentbox=congaree, below=0.8cm of llm2] (term2) {Terminate};
\draw[dataarrow, dashed] (llm2) -- (term2);

% Main flow
\draw[dataarrow] (llm2) -- (output);
\draw[dataarrow] (output) -- (parse);
\draw[dataarrow] (parse) -- (t1);
\draw[dataarrow] (t1) -- (t2);
\draw[dataarrow] (t2) -- (t3);
\draw[dataarrow] (t3) -- (allres);
\draw[dataarrow] (allres.north) -- ++(0,0.8) -| (llm2.north);

\end{tikzpicture}
\caption{Batch execution workflow. The LLM generates a complete response, tool calls are parsed and executed serially (T1, T2, T3), and all results are fed back to trigger the next reasoning cycle.}
\label{fig:batch}
\end{figure*}

\begin{figure*}[t]
\centering
\begin{tikzpicture}[node distance=1.2cm and 1.8cm]

% Parallel tasks
\node[pbar=atlantic, minimum width=2.5cm] (t1) {Task 1 \quad $\parallel$};
\node[pbar=atlantic, minimum width=2.5cm, below=0.6cm of t1] (t2) {Task 2 \quad $\parallel$};

% Serial tasks
\node[pbar=horseshoe, minimum width=2.5cm, right=2.5cm of $(t1)!0.5!(t2)$] (t3) {Task 3};
\node[pbar=horseshoe, minimum width=2.5cm, right=1.5cm of t3] (t4) {Task 4};

% Arrows from parallel into serial
\draw[dataarrow] (t1.east) -- (t3.west);
\draw[dataarrow] (t2.east) -- (t3.west);
\draw[dataarrow] (t3) -- (t4);

% Annotations
\draw[thick, decorate, decoration={brace, amplitude=5pt, mirror, raise=3pt}]
  (t1.north west) -- (t2.south west)
  node[midway, left=0.4cm, font=\footnotesize, align=center] {concurrent};
\draw[thick, decorate, decoration={brace, amplitude=5pt, raise=3pt}]
  (t3.north west) -- (t4.north east)
  node[midway, above=0.3cm, font=\footnotesize, align=center] {sequential};

\end{tikzpicture}
\caption{Parallel execution: tasks marked $\parallel$ run concurrently. Once both complete, subsequent tasks execute sequentially. The agent decides which tasks are independent and can safely overlap.}
\label{fig:parallel_workflow}
\end{figure*}
